\documentclass[10pt,a4paper]{amsart}
\usepackage[margin=19mm]{geometry}
\usepackage{amsmath,amssymb,mathtools}
\usepackage[scr=rsfs,cal=euler]{mathalfa}
\usepackage{xfrac}
\usepackage[unicode,hidelinks,bookmarks=false]{hyperref}
\newcommand{\ie}{i.e.\ }
\newcommand{\cf}{cf.\ }
\newcommand{\resp}{respectively\ }
\newcommand{\Subsection}{\subsection}
\providecommand{\nobreakdash}{\nobreak\hbox{-}}
\setlength{\emergencystretch}{2em}
\multlinegap0pt
\usepackage{bm}
\usepackage{bbm}
\usepackage{dsfont}
\usepackage{xspace}
\usepackage{cancel}
\usepackage{mathtools}
\usepackage{amsthm}
\makeatletter
\@ifundefined{theorem}{\newtheorem{theorem}{Theorem}}{}
\@ifundefined{lemma}{\newtheorem{lemma}[theorem]{Lemma}}{}
\makeatother
\newcommand{\CR}{\mathcal{R}}
\newcommand{\Dphi}{\Dot{\phi}}
\newcommand{\Drho}{\Dot{\rho}}
\newcommand{\im}[1]{\mathrm{Im} \left(#1\right)}
\newcommand{\re}[1]{\mathrm{Re}\left(#1\right)}
\title[On the Diverse Dynamical Behaviors Arising in Deep Linear Tr]{On the Diverse Dynamical Behaviors Arising in Deep Linear Transformers}
\author{Sixu Li, Thomas Jacob Maranzatto, Jan Peszek, Trevor Teolis, Semih Akkoc, Konstantin Riedl, Sennur Ulukus, Nicol\'as Garc\'ia Trillos}
\date{Source: arXiv:2607.18584v1 (CC BY 4.0)}
\begin{document}
\maketitle

\noindent\emph{Source and licence.} On the Diverse Dynamical Behaviors Arising in Deep Linear Transformers. Version arXiv:2607.18584v1 (CC BY 4.0), distributed under the Creative Commons Attribution 4.0 International licence, as recorded in the arXiv licence metadata for this exact version and shown on the abstract page https://arxiv.org/abs/2607.18584v1. Full rights notice: licenses/arxiv-2607.18584-CC-BY-4.0.txt.\par\smallskip

\noindent\textit{Editorial scope.} Counted unit: the Lemma whose source label is \texttt{lem:pc\_dyn\_general\_expand} in the article (source0.tex lines 5606--5616, source label \texttt{lem:pc\_dyn\_general\_expand}), delivered together with the article's own complete original proof of it (source0.tex lines 5618--5755, one \texttt{proof} environment of 138 lines). No mathematics is written, repaired or re-derived: statement and proof are the article's own text line for line. Required context retained uncounted: eq:dyn\_r2 (lines 1188--1191); eq:pc\_dyn\_general (lines 1250--1252); eq:coeff\_b (lines 3342--3347); eq:coeff\_B\_C (lines 3349--3358). Every definition, display and label of those lines is retained unchanged. Omitted from this package and not redistributed: 1--1187, 1192--1249, 1253--3341, 3348--3348, 3359--5605, 5617--5617, 5756--6943. Nothing inside a retained excerpt is omitted or altered: every retained body is delivered exactly as the article's lines are written, so no omission receipt of any kind accompanies this package. Citations: the retained excerpts carry no citation command at all, so no bibliography entry of the article is transcribed and no reference is redistributed. Reference adaptation: none was needed and none was made; every reference inside the delivered unit resolves inside it, so no unresolved reference or citation is produced. Printed slips kept literal: no printed word, symbol, hypothesis or number of the counted statement or of its proof was altered. Layout and class: the article's own class is \texttt{article} with packages including inputenc, footmisc, amsmath, graphicx, amsthm, tikz-cd, amssymb, bm, bbm, booktabs, pifont, xcolor, array, makecell; the wrapper is the lane's pinned \texttt{amsart} editorial wrapper at 10pt on A4, which restates the theorem-like environments the retained text needs and adds the article's own macro definitions that the retained text needs (\texttt{\textbackslash CR}, \texttt{\textbackslash Dphi}, \texttt{\textbackslash Drho}, \texttt{\textbackslash im}, \texttt{\textbackslash re}), with extra wrapper packages (bm, bbm, dsfont, xspace, cancel, mathtools). The delivered numbering is the unit's own. Assets: no retained span contains a figure environment, a graphics-inclusion command, a PDF-page inclusion, a table or a TikZ picture; no image file is copied into this package, the package declares no asset dependency and no image file is redistributed with it. Licence: version arXiv:2607.18584v1 of this article is distributed under the Creative Commons Attribution 4.0 International licence, as recorded in the arXiv licence metadata for this exact version and shown on the abstract page https://arxiv.org/abs/2607.18584v1. A full rights notice is delivered at licenses/arxiv-2607.18584-CC-BY-4.0.txt. Characters: every retained excerpt is pure ASCII, so the delivered engine, XeLaTeX with its default Latin Modern text font, reports no missing character.
    \begin{equation}\label{eq:dyn_r2}
    \Dot{\CR}_2(t) = F(\CR_2(t), \CR_2(t))
    = 2i \left( B(\CR_2(t)) (\CR_2(t))^2 + C(\CR_2(t)) \CR_2(t) + \overline{B}(\CR_2(t)) \right) .
    \end{equation}

\begin{equation}\label{eq:pc_dyn_general}
\Drho(t) = \re{e^{-i\phi(t)} F(\rho(t) e^{i\phi(t)})}, \qquad \Dphi(t) = \frac{1}{\rho(t)} \im{e^{-i\phi(t)} F(\rho(t) e^{i\phi(t)})} .
\end{equation}

\begin{equation}\label{eq:coeff_b}
    \begin{gathered}
    b_1 := \frac{1}{16}(iw_5+w_6+w_7-iw_8), \qquad b_2 := \frac{1}{16}(iw_5-w_6+w_7+iw_8), \qquad b_3 := \frac{1}{8}(iw_1-w_2),\\
    b_4 := \frac{1}{8}(-w_3+iw_4), \qquad b_5 := -\frac{1}{4}w_0.
    \end{gathered}
\end{equation}

\begin{equation}
    \label{eq:coeff_B_C}
         B(r)
        := b_1 r + b_2 \bar{r} + b_3
        \qquad
        \text{and}
        \qquad
        C(r)
        := b_4 r + \overline{b_4} \bar{r} + b_5.
    \end{equation}

\begin{lemma}\label{lem:pc_dyn_general_expand}
The polar-coordinate system \eqref{eq:pc_dyn_general} admits the equivalent representation
\begin{equation*}
    \begin{split}
        \Drho = 2(1-\rho^2) \left( S_1(\phi) \rho + S_2(\phi) \right), \qquad
        \Dphi = 2 \left( S_3(\phi) (\rho^2+1) + S_4(\phi) \frac{\rho^2 + 1}{\rho} + 2S_5(\phi) \rho + \re{b_5} \right),
    \end{split}
    \end{equation*}
where both functions $S_{j}$ constant $b_5$ depend on the matrices $A$ and $V$.
The explicit form of $S_j$ and $b_5$ are provided in \eqref{eq:func_Sj} and \eqref{eq:coeff_b}, respectively.
\end{lemma}

\begin{proof}
    For notational convenience we drop the time dependence on all variables.
    
    Substituting $\CR_2 = \rho e^{i\phi}$ and $\overline{\CR_2} = \rho e^{-i\phi}$ into the ODE \eqref{eq:dyn_r2} satisfied by $\CR_2$ gives
    \begin{equation}\label{eq:pc_transform}
        \frac{d}{dt} \CR_2 = \frac{d}{dt} (\rho e^{i\phi}) = e^{i\phi} \left( \Drho + i \rho \Dphi \right) .
    \end{equation}
    By using the definitions of $B$ and $C$ from \eqref{eq:coeff_B_C}, we can expand the left-hand side of the above equation:
    \begin{equation*}
    \begin{aligned}
        \mathrm{LHS} &= \frac{d}{dt} \CR_2 = 2i \left( B(\CR_2) (\CR_2)^2 + C(\CR_2) \CR_2 + \overline{B}(\CR_2) \right) \\
        &= 2i \left( b_1 \CR_2^3 + b_2 \CR_2^2\overline{\CR_2} + b_3\CR_2^2 + b_4 \CR_2^2 + \overline{b_4} \overline{\CR_2}\CR_2 + b_5 \CR_2 + \overline{b_1} \overline{\CR_2} + \overline{b_2} \CR_2 + \overline{b_3} \right) \\
        &= 2i \left( b_1 \rho^3 e^{i3\phi} + b_2 \rho^3 e^{i\phi} + (b_3 + b_4) \rho^2 e^{i2\phi} + \overline{b_4} \rho^2 + (\overline{b_2} + b_5) \rho e^{i\phi} + \overline{b_1} \rho e^{-i\phi} + \overline{b_3} \right)\\
        &= e^{i\phi} \left(2 \left( ib_1e^{i2\phi} \rho^3 + i \overline{b_1} e^{-i2\phi} \rho + i(b_3 + b_4) e^{i\phi} \rho^2 + i (\overline{b_3} + \overline{b_4} \rho^2) e^{-i\phi} + i \left( b_2 \rho^3 + (\overline{b_2} + b_5) \rho \right) \right) \right) = e^{i\phi} G(\rho, \phi),
    \end{aligned}
    \end{equation*}
    where we defined $G$ implicitly in the last step.
   
    Comparing the two expressions for the left-hand and right-hand sides of~\eqref{eq:pc_transform} results in the complex-valued identity
    \begin{equation}
    \label{eq:complex_iden}
        \Drho + i \rho \Dphi = G(\rho, \phi).
    \end{equation}
    Recalling that $\rho$ and $\phi$ are both real-valued,
    we separate the real and imaginary parts of the function $G$.
    In particular, letting
    \begin{equation*}
        G(\rho, \phi)
        = G_R (\rho, \phi) + i G_I (\rho, \phi)
        \qquad
        \text{with}
        \qquad
        G_R := \re{G}
        \text{ and }
        G_I := \im{G},
    \end{equation*}
    we can match the real and imaginary parts of \eqref{eq:complex_iden} to obtain the real-valued coupled ODE system
    \begin{equation}
        \label{eq:aux_coupled_dyn}
        \Drho = G_R(\rho, \phi), \qquad
        \Dphi = \frac{1}{\rho} G_I(\rho, \phi),
    \end{equation}
    and it remains to write down $G_R$ as well as $G_I$ explicitly.
    
    To do so, we introduce some notation and set
    \begin{equation*}
    \begin{aligned}
        G(\rho, \phi)
        &= 2 \left( ib_1e^{i2\phi} \rho^3 + i \overline{b_1} e^{-i2\phi} \rho + i(b_3 + b_4) e^{i\phi} \rho^2 + i (\overline{b_3} + \overline{b_4} \rho^2) e^{-i\phi} + i \left( b_2 \rho^3 + (\overline{b_2} + b_5) \rho \right) \right)\\
        &= 2 \left( A_1 + A_2 + A_3 + A_4 + A_5 \right).
    \end{aligned}
    \end{equation*}
    % For all $j=1,2,\dots, 5$,
    % we denote by $a_j$ and $c_j$ the real and imaginary parts of $b_j$, i.e., $b_j = a_j + i c_j$.
    % Moreover, 
    % using that we have for any $\ell \in \mathbb{N}_+$, and $a, c \in \R$ the identities
    % \begin{equation*}
    %     \re{(a + ic)e^{i\ell\phi}} = a \cos (\ell\phi) - c\sin (l\phi)
    %     \qquad
    %     \text{and}
    %     \qquad
    %     \im{(a + i c) e^{i\ell\phi}}
    %     = c \cos (l\phi) + a \sin (\ell\phi), 
    % \end{equation*}
    We can then compute for each term $A_j$ the real and imaginary part.
    \begin{itemize}
        \item Term $A_1$:
        \begin{equation*}
        \begin{aligned}
        \re{A_1} = \re{ib_1 e^{i2\phi} \rho^3} = - \im{b_1 e^{i2\phi}} \rho^3, \qquad 
        \im{A_1} = \im{ib_1 e^{i2\phi} \rho^3} = \re{b_1 e^{i2\phi}} \rho^3.
        \end{aligned}
        \end{equation*}
    
        \item Term $A_2$:
        \begin{equation*}
        \begin{aligned}
        \re{A_2} = \re{i\overline{b_1} e^{-i2\phi} \rho} = \im{b_1 e^{i2\phi}} \rho, \qquad \im{A_2} = \im{i\overline{b_1} e^{-i2\phi} \rho} = \re{b_1 e^{i2\phi}} \rho .
        \end{aligned}
        \end{equation*}
    
        \item Term $A_3$:
        \begin{equation*}
        \begin{aligned}
        &\re{A_3} = \re{i(b_3 + b_4) e^{i\phi} \rho^2} = - \im{(b_3 + b_4) e^{i\phi}} \rho^2,\\[5pt]
        &\im{A_3} = \im{i(b_3 + b_4) e^{i\phi} \rho^2} = \re{(b_3 + b_4) e^{i\phi}} \rho^2 .
        \end{aligned}
        \end{equation*}
    
        \item Term $A_4$:
        \begin{equation*}
        \begin{aligned}
        &\re{A_4} = \re{i (\overline{b_3} + \overline{b_4} \rho^2) e^{-i\phi}} = \im{b_3 e^{i\phi}} + \im{b_4 e^{i\phi}} \rho^2,\\[5pt]
        &\im{A_4} = \im{i (\overline{b_3} + \overline{b_4} \rho^2) e^{-i\phi}} = \re{b_3 e^{i\phi}} + \re{b_4 e^{i\phi}} \rho^2 .
        \end{aligned}
        \end{equation*}
    
        \item Term $A_5$:
        \begin{equation*}
        \begin{aligned}
        &\re{A_5} = \re{i \left( b_2 \rho^3 + (\overline{b_2} + b_5) \rho \right)} = -\im{b_2} \rho^3 + \im{b_2} \rho,\\[5pt]
        &\im{A_5} = \im{i \left( b_2 \rho^3 + (\overline{b_2} + b_5) \rho \right)} = \re{b_2} \rho^3 +  (\re{b_2} + \re{b_5}) \rho .
        \end{aligned}
        \end{equation*}
    \end{itemize}
    For the last term, recall that $b_5$ is real-valued.
    
    With those equalities we have 
    \begin{equation*}
    \begin{aligned}
    G_R(\rho, \phi) &= 2 \left( \re{A_1} + \re{A_2} + \re{A_3} + \re{A_4} + \re{A_5} \right)\\
    &= 2 (1-\rho^2)\left( \im{b_1 e^{i2\phi} + b_2} \rho  + \im{b_3 e^{i\phi}} \right),
    \end{aligned}
    \end{equation*}
    and
    \begin{equation*}
    \begin{aligned}
    G_I(\rho, \phi) &= 2 \left( \im{A_1} + \im{A_2} + \im{A_3} + \im{A_4} + \im{A_5} \right)\\
    &= 2 \left( (\rho^2 + 1) \left(  \re{b_1 e^{i2\phi} + b_2} \rho + \re{b_3 e^{i\phi}} \right) + \rho \left(2 \re{b_4 e^{i\phi}} \rho + \re{b_5} \right) \right) .
    \end{aligned}
    \end{equation*}
    By introducing the notation
    \begin{equation}
    \label{eq:func_Sj}
    \begin{aligned}
    &\qquad \qquad S_1(\phi) := \im{b_1 e^{i2\phi} + b_2}, \qquad S_2(\phi) := \im{b_3 e^{i\phi}},\\[5pt]
    &S_3(\phi) := \re{b_1 e^{i2\phi} + b_2}, \qquad S_4(\phi) := \re{b_3 e^{i\phi}}, \qquad S_5(\phi) := \re{b_4 e^{i\phi}},
    \end{aligned}
    \end{equation}
    we eventually obtain the coupled ODE system
    \begin{equation*}
    \begin{split}
        \Drho &= 2 (1-\rho^2) \left( S_1(\phi) \rho + S_2(\phi) \right)\\
        \Dphi &= 2 \left( S_3(\phi) (\rho^2+1) +  S_4(\phi) \frac{\rho^2 + 1}{\rho} + 2 S_5(\phi) \rho + \re{b_5} \right),
    \end{split}
    \end{equation*}
    which concludes the proof.
\end{proof}

\end{document}
